% ARCHIVED DRAFT: independently audited by root and incorporated on 2026-11-19 16:42 UTC.
% The assembled paper is authoritative; this file retains drafting provenance.
\begin{saresult}{B.071}{Proved}{All finite-residue local bases with principal maximal ideal}
Let $(A,\mathfrak m)$ be any nonzero local commutative ring with finite
residue field and $\mathfrak m=\pi A$, and let $A\subset S$ be a unital
commutative inclusion. Put
\[
                  C=\{s\in S:\pi^Ns\in A
                                  \text{ for some }N\ge0\}.
\]
Then the polynomial coaddition and comultiplication tensors are uniquely
determined precisely in the following cases:
\[
 C=A,\quad\text{or}\quad
 \pi\text{ is regular in }A\text{ and }C=A[1/\pi]
 \text{ as a canonical subring of }S.                     \tag{B71.1}
\]
Every positive case has all polynomial tensor isomorphisms.
If $\pi$ is regular in $A$, the criterion is equivalently
$\operatorname{Ann}_S(\pi)=0$.
If $\pi$ is nilpotent, the criterion reduces to $S=A$.
Neither Noetherianity nor reducedness is assumed.
\end{saresult}
\begin{proof}
If $\pi$ is nilpotent, the finite residue field and the successive cyclic
quotients $\pi^jA/\pi^{j+1}A$ make $A$ finite. Also $C=S$, and B.013
gives exactly the stated conclusion. This includes the field case
$\pi=0$. We henceforth assume $\pi$ is not nilpotent.

Set $P=\bigcap_{j\ge0}\pi^jA$. First
\[
          \operatorname{Ann}_A(\pi)\subset P,
          \qquad P=\pi P,\qquad \bar A=A/P
          \text{ is a discrete valuation ring}.           \tag{B71.2}
\]
If $x\notin P$, choose the greatest $j$ with $x\in\pi^jA$ and write
$x=\pi^ju$; then $u\notin\mathfrak m$, so $u$ is a unit.
Such an $x$ cannot be annihilated by $\pi$, as that would make
$\pi^{j+1}=0$. This proves the first inclusion without cancelling a
zero divisor. For $x\in P$ write $x=\pi y$. For every $j$ one can also
write $x=\pi^{j+1}z_j$, so
$y-\pi^jz_j\in\operatorname{Ann}_A(\pi)\subset P$.
Hence $y\in P$, proving $P=\pi P$.
If $\pi x\in P$, the same equations show $x\in P$, so $\pi$ is regular
in $\bar A$. Every nonzero element of $\bar A$ is a unit times a finite
power of $\pi$. Products are nonzero and every nonzero ideal is generated
by an element of smallest exponent. Thus $\bar A$ is a DVR, as claimed.

All coefficients of integer-valued polynomials belong to $C$.
Indeed the quotients $A/\pi^nA=\bar A/\pi^n\bar A$ are finite and have
the usual discrete-valuation digit representatives. Choose enough
distinct such representatives for a Vandermonde grid. Every difference
is a unit times a finite power of $\pi$, so Cramer's rule gives the
asserted saturation for each coefficient. Conversely, if $\pi^Nc\in A$,
take the monic polynomial that is the product of $X-a$ over a complete
set of representatives of $A/\pi^NA$. Every value is in $\pi^NA$, so
$c$ occurs as its multiplied leading coefficient. Thus $C$ is the actual
coefficient ring and $\Omega(A,S)=\Omega(A,C)$.

One has $PC=P$. For $x\in P$ and $\pi^Nc=a\in A$, write
$x=\pi^{N+j}b_j$ for every $j$, giving $xc=\pi^jb_ja\in P$.
It follows, in every number of variables, that reduction gives
\[
  \Omega_n(A,C)/P\Omega_n(A,C)
       \simeq\Omega_n(\bar A,\bar C),\qquad
  P\Omega_n(A,C)=P[X_1,\ldots,X_n],\quad \bar C=C/P.       \tag{B71.3}
\]
Every reduced polynomial lifts because its lifted values lie in the
preimage of $\bar A$, namely $A$; the asserted kernel follows from
$PC=P$.

Crucially, tensor maps lift as well as descend. Let $M=\Omega(A,C)$.
The restriction of the polynomial tensor map is an isomorphism
\[
              P M^{\otimes_A n}\ \simeq P[X_1,\ldots,X_n]. \tag{B71.4}
\]
Here is an elementwise proof, valid with zero divisors. For $h\in P$
and $F\in M$, choose $N$ with $\pi^NF\in A[X]$ and write
$h=\pi^Nh'$ with $h'\in P$, using (B71.2). Moving $\pi^N$ across a tensor
sign expands a factor $h\otimes F$ as the ordinary monomial tensor with
coefficients $h$ times the coefficients of $F$, all of which belong to
$P$ because $PC=P$. For two factors the explicit equality is
$hF\otimes G=h'F\otimes\pi^NG$ when $h=\pi^Nh'$ and
$\pi^NG\in A[X]$; expand the second factor and use $PC=P$
for every coefficient of the first.
Repeating this step in each factor writes every
element of $PM^{\otimes n}$ as that ordinary monomial tensor of its
polynomial image. Conversely these monomial tensors exist, and their
polynomial images recover their coefficients. This proves both
injectivity and surjectivity in (B71.4).

Reduction of $M^{\otimes n}$ modulo $P$ is the tensor power over
$\bar A$ of the reduced module. Together with (B71.3)--(B71.4), this
shows that every $\theta_n$ is injective, surjective, or an isomorphism
exactly when its reduced counterpart is. For injectivity, reduce a
kernel element and then use (B71.4); for surjectivity, lift a reduced
preimage and correct the remaining polynomial in $P[X_1,\ldots,X_n]$
using (B71.4). The same correction proves that coaddition or
comultiplication exists upstairs exactly when it exists downstairs.
These arguments require no flatness of $A/P$.

Apply the DVR case B.063 to $\bar A\subset\bar C$. Every coefficient
of $\bar C$ is $\pi$-power saturated over $\bar A$, so the positive cases
are exactly $\bar C=\bar A$ and $\bar C=\operatorname{Frac}(\bar A)$.
The first says $C=A$. In the second, lift an inverse of $\pi$ to obtain
$\pi c=1+x$ with $x\in P\subset\mathfrak m$. The element $1+x$ is a unit
of $A$, so $\pi$ is a unit of $C$. Injectivity of $A\subset C$ therefore
forces $\pi$ to be regular in $A$, and the definition of $C$ gives
$C=A[1/\pi]$. Both cases lift with all tensor isomorphisms by
(B71.3)--(B71.4).

Finally suppose $\pi$ is regular in $A$. If $\pi$ is regular in $S$,
then $C\subset A[1/\pi]$; any $a/\pi^N\notin A$ has
$a=\pi^ju$ for $j<N$ and $u$ a unit, so it forces $\pi^{-1}\in C$.
Thus (B71.1) holds. Conversely an element annihilated by $\pi$ lies in
$C$, and either positive option in (B71.1) has no such nonzero element.
This proves the equivalent annihilator criterion.
\end{proof}
